\subsection{3. Quantitative Relationship between Polarization Differences and the Zone-Averaged Shift Photocurrent}
\textbf{OSTI:} \texttt{1523841}\quad\textbf{Authors:} Fregoso, Morimoto, Moore (2017)\quad\textbf{Domain:} Condensed Matter / Topological Response / Shift Current Photovoltaics\\
\scorebadge{Coverage}{7}\quad\scorebadge{Agreement}{10}\quad\textbf{Total:} 17/20\par\smallskip
\noindent\textbf{Coverage rationale.} Reproduced the paper's central analytical identity e*R\_bar\_cv = a*(P\_c-P\_v) + W\_cv*e*a (Eq. 9) on the Rice-Mele model (paper's Appendix D testbed). Verified d-vector shift-vector formula (Eq. C5), analytic limit formulae (Eqs. D8-D9), Berry-connection closed forms (Eq. D4), and the winding-number discontinuity at the gap-closing point delta=0. NOT replicated: the general multi-band formula (Section III/Eq. 17), material-specific DFT+Wannier calculations (GeS, BaTiO3, etc.), nor the frequency-integrated shift-conductivity connection.\par
\noindent\textbf{Agreement rationale.} Identity verified to machine precision: residual |e*R\_bar\_cv - a*(P\_c-P\_v)| \textasciitilde{} 3.3e-16 pointwise on a 100-point delta sweep. Analytic limit formulas reproduced to \textasciitilde{}1e-16. Winding number |W\_cv|=1 correctly identified with Delta\_W=1 jump across delta=0. This is essentially a perfect analytical replication on the 1D two-band model the paper uses as its primary illustration.\par\smallskip
\noindent\textbf{What reproduced:}
\begin{itemize}[leftmargin=1.4em,itemsep=0.2em,topsep=0.2em]
\item Paper Eq. 9: e*R\_bar\_cv = a(P\_c-P\_v) + W\_cv*e*a
\item Paper Eq. C5: d-vector shift-vector closed form
\item Paper Eqs. D8-D9: analytic k-\textgreater{}0 and k-\textgreater{}pi/a limits
\item Paper Eq. D4: Berry connections A\_nn(k)
\item Winding-number integer discontinuity at gap closing
\item Fig. 1(b) qualitative structure: P\_c-P\_v vs d-vec vs full numerical R\_bar
\item Gauge-unique closed-form Berry-connection integration
\end{itemize}\noindent\textbf{What missing:}
\begin{itemize}[leftmargin=1.4em,itemsep=0.2em,topsep=0.2em]
\item General N-band formula (Section III / Eq. 17)
\item Material-specific first-principles + Wannier calculations
\item Frequency-integrated shift-conductivity connection
\item Finite-frequency shift current spectra sigma\textasciicircum{}shift(omega)
\item Paper's Figs. 2-4 (material examples, higher-D extensions)
\item Comparison against experimental shift-current measurements
\end{itemize}\noindent\textbf{Follow-on research questions:}
\begin{enumerate}[leftmargin=1.6em,itemsep=0.2em,topsep=0.2em,label=Q\arabic*.]
\item Extend the numerical verification to a 3-band or 4-band tight-binding model (e.g. SSH+onsite, or a minimal GeS-like model) and check whether the identity generalizes via the Wilson-loop formulation of Eq. 17.
\item Compute the shift-current conductivity sigma\textasciicircum{}shift(omega) in the Rice-Mele model and verify the frequency integral equals (pi*e\textasciicircum{}3/hbar\textasciicircum{}2)*(P\_c - P\_v) modulo the winding contribution.
\item Apply the identity as a DFT+Wannier post-processing tool for a material like GeS or monolayer WS2: compute P\_c-P\_v from Wannier centers and compare to the direct frequency-integrated sigma\textasciicircum{}shift.
\item Study how the winding number W\_cv changes under SHG-relevant perturbations (strain, electric field) and predict a shift-current discontinuity as a material tuning knob.
\item Generalize to interacting systems: does the identity survive a Hartree-Fock or GW correction when P\_c-P\_v is evaluated with quasiparticle wavefunctions Test on a Hubbard-Rice-Mele numerical example.
\end{enumerate}

